\begin{abstract}
Large vision-language models answer visual questions from long visual-token sequences, and that prefix dominates prefill compute and the key-value cache. Existing pruners fall into three families, and each family leaves a distinct gap. Salience ranking keeps high-scoring tokens that still contain a large amount of similar content, and a score from one visual depth cannot represent what other depths emphasize. Diversity and coverage build a compact visual set, but image-level representativeness and the evidence required by the current question are not fully consistent. One-shot compression before the language model then freezes the visual length, so early interaction and later redundancy share one capacity. We formulate pruning as query-adaptive subset selection under a layered compute budget and a coverage-fidelity constraint: near a fixed-capacity visual basis, raise question relevance while limiting coverage distortion and the number of identity edits. The training-free framework Q-SCOPE builds that basis from two visual depths and regional quotas, adapts it with query-guided exchange, and contracts the nested set with question-to-visual attention. On LLaVA-1.5-7B, the completed early schedule removes 94.4\% of the visual token-layer budget and retains 95.66\% of the unpruned six-task relative mean. Cells without a matched measurement are left blank. The relative mean is not a per-question accuracy, and the budget ratio is not a measured speedup.
\end{abstract}
